\documentclass[11pt,a4paper]{article}
\usepackage{fontspec,amsmath,amssymb,amsthm,microtype}
\setmainfont{lmroman10-regular.otf}[BoldFont=lmroman10-bold.otf,ItalicFont=lmroman10-italic.otf,BoldItalicFont=lmroman10-bolditalic.otf]
\setmonofont{lmmono10-regular.otf}[Scale=MatchLowercase]
\usepackage[margin=25mm]{geometry}
\usepackage[unicode,hidelinks]{hyperref}
\setlength{\emergencystretch}{3em}
\setlength{\parskip}{0.4em}
\setlength{\parindent}{0pt}
\allowdisplaybreaks[2]
\raggedbottom
\title{\texorpdfstring{A concrete latent/\allowbreak{}full Gaussian transition coupling}{A concrete latent/full Gaussian transition coupling}}
\author{LangevinLean mathematical exposition}
\date{}
\begin{document}\maketitle
\noindent Source module: \texttt{Proximal\allowbreak{}Langevin/\allowbreak{}MYPGDTransport.\allowbreak{}lean}. The accompanying HTML pairs this mathematical exposition with the exact Lean statements and records the scope of the formalization.


For an \(r\)-Lipschitz deterministic map \(D\), compare \(D(x)+aP_x\xi\) and \(D(y)+a\xi\) using the same fresh standard Gaussian vector. Initial laws are probability measures with finite second moments.


\section{\texorpdfstring{Lipschitz maps and explicit pairs have finite transport cost}{Lipschitz maps and explicit pairs have finite transport cost}}


Center \(D\) at zero to use \(\|D(x)-D(0)\|\le r\|x\|\), then add back its constant value; this proves preservation of \(L^2\). For two \(L^2\) variables on one space, their joint law is a coupling. Its cost bounds \(W_2\), yielding both finiteness and \(W_2(\mathcal L(X),\mathcal L(Y))^2\le\mathbb E\|X-Y\|^2\).




A globally Lipschitz map \(D\) need not vanish at zero. The first lemma writes \(D(X)=[D(X)-D(0)]+D(0)\): the centered map preserves \(L^2\), and the constant belongs to \(L^2\) because the measure is a probability. For the transport lemma, convert \(X,Y\) to their \(L^2\) classes. Their joint realization is an admissible coupling of their pushforwards, so \[W_2(X_\#P,Y_\#P)^2\le\|X-Y\|_{L^2(P)}^2=\int\|X-Y\|^2dP.\] The representative equalities \texttt{hx}, \texttt{hy} prove that these pushforwards are the original laws, and the finite norm bound proves finite Wasserstein distance before squaring its real value.




\textbf{Correspondence with the displayed declarations.}


\begin{itemize}

\item \texttt{mem\allowbreak{}Lp\_\allowbreak{}lipschitz\_\allowbreak{}comp}: Globally Lipschitz maps preserve finite quadratic moments on a probability space.
\item \texttt{wasserstein\_\allowbreak{}sq\_\allowbreak{}le\_\allowbreak{}integral\_\allowbreak{}pair}: Direct finite-cost transport estimate for an actual L\ensuremath{{}^{2}} pair.
\end{itemize}


{\raggedright\noindent Lean declarations: \texttt{mem\allowbreak{}Lp\_\allowbreak{}lipschitz\_\allowbreak{}comp}, \texttt{wasserstein\_\allowbreak{}sq\_\allowbreak{}le\_\allowbreak{}integral\_\allowbreak{}pair}.\par}

\section{\texorpdfstring{Measurable Gaussian maps preserve second moments}{Measurable Gaussian maps preserve second moments}}


The two maps are \(T_x(u,\xi)=D(u)+aP_x\xi\) and \(T(u,\xi)=D(u)+a\xi\). Lipschitz continuity of \(D\), linearity of \(P_x\), addition and scalar multiplication prove joint measurability. On \(\mu\otimes\mathcal N(0,I)\), each summand belongs to \(L^2\), so both pushforward transition laws have finite second moments.




Both updates use a fresh \(Z\sim\mathcal N(0,I)\), independent of the input: \(D(x)+aP_xZ\) and \(D(x)+aZ\). Their measurability follows by composing continuous \(D\), the coordinate projections, addition and scalar multiplication. For moment propagation, \texttt{comp\_\allowbreak{}fst} pulls the \(L^2\) drift onto the input–Gaussian product and \texttt{comp\_\allowbreak{}snd} pulls the Gaussian noise onto that same space. The sum is in \(L^2\); \texttt{mem\allowbreak{}Lp\_\allowbreak{}map\_\allowbreak{}measure\_\allowbreak{}iff} then transfers this fact to the output law. Thus the next coupling argument has actual finite-moment marginals.




\textbf{Correspondence with the displayed declarations.}


\begin{itemize}

\item \texttt{latent\allowbreak{}Gaussian\allowbreak{}Update}: The deterministic map followed by latent-only Gaussian noise.
\item \texttt{full\allowbreak{}Gaussian\allowbreak{}Update}: The same map followed by full Gaussian noise.
\end{itemize}


{\raggedright\noindent Lean declarations: \texttt{latent\allowbreak{}Gaussian\allowbreak{}Update}, \texttt{full\allowbreak{}Gaussian\allowbreak{}Update}, \texttt{latent\allowbreak{}Gaussian\allowbreak{}Update\_\allowbreak{}measurable}, \texttt{full\allowbreak{}Gaussian\allowbreak{}Update\_\allowbreak{}measurable}, \texttt{latent\allowbreak{}Gaussian\allowbreak{}Step\_\allowbreak{}mem\allowbreak{}Lp}, \texttt{full\allowbreak{}Gaussian\allowbreak{}Step\_\allowbreak{}mem\allowbreak{}Lp}.\par}

\section{\texorpdfstring{Extend every input coupling by an independent Gaussian}{Extend every input coupling by an independent Gaussian}}


For any input coupling \(c\) of \(\mu,\nu\), work on \(c\otimes\mathcal N(0,I)\). The two outputs have precisely the required Markov marginals, and their difference is \(D(x)-D(y)-aP_\theta\xi\). The exact removed-noise identity gives squared coupling cost \(\mathbb E_c\|D(x)-D(y)\|^2+a^2d_\theta\), bounded by \(r^2\operatorname{cost}(c)^2+a^2d_\theta\). This constructs the recurrence input, with no optimal-coupling assumption.




Start from any input coupling \(c\) of \(\mu,\nu\), and work on \(c\otimes\mathcal N(0,I)\). Define \(F=D(X)+aP_xZ\) and \(H=D(Y)+aZ\). The facts \texttt{h\allowbreak{}Flaw} and \texttt{h\allowbreak{}Hlaw} check separately that their marginals are the desired two Markov transitions; this ensures the constructed pair is a coupling of the actual algorithms.



The cancellation in \texttt{hid} is \[F-H=D(X)-D(Y)-aP_\theta Z.\] By the product Gaussian identity, its mean square is \(\mathbb E_c\|D(X)-D(Y)\|^2+a^2d_\theta\). If \(D\) is \(r\)-Lipschitz, the first term is bounded by \(r^2\mathbb E_c\|X-Y\|^2\). Finally \texttt{hnorm} identifies the latter integral with the square of the coupling cost. All three random variables are proved square integrable before any of these integral identities is used.




\textbf{Correspondence with the displayed declarations.}


\begin{itemize}

\item \texttt{latent\_\allowbreak{}full\_\allowbreak{}markov\_\allowbreak{}coupling\_\allowbreak{}sq}: Actual latent-noise and full-noise one-step laws admit the stated squared coupling bound. The independent centered parameter noise supplies exactly \texttt{a\ensuremath{{}^{2}}*dtheta}, and has no mixed contribution.
\end{itemize}


{\raggedright\noindent Lean declarations: \texttt{latent\_\allowbreak{}full\_\allowbreak{}markov\_\allowbreak{}coupling\_\allowbreak{}sq}.\par}

\end{document}
